\documentclass[11pt,a4paper]{article}
\usepackage[margin=25mm]{geometry}
\usepackage{amsmath,amssymb,amsthm}
\usepackage{booktabs}
\usepackage{hyperref}
\usepackage{enumitem}
\setlist{nosep}

\newtheorem{proposition}{Proposition}
\newtheorem{lemma}{Lemma}
\newtheorem{definition}{Definition}
\newcommand{\N}{\mathcal N}
\newcommand{\E}{\mathbb E}
\newcommand{\R}{\mathbb R}
\newcommand{\Cov}{\operatorname{Cov}}
\newcommand{\pos}[1]{(#1)_+}

\title{Hidden hills in the old basket model,\\ and why the mixture form has none}
\author{Ajit}
\date{28 October 2026}

\begin{document}
\maketitle

\begin{abstract}
Both the old Version-4 basket model and its mixture form are built from the same exact law of a basket
given a latent taste $z$. They differ only in the weight each taste receives. In the old model that weight
is the Gaussian density times $Z_z$, the total score of all baskets the taste allows. Because $Z_z$ grows
fastest for tastes that make large baskets attractive, the weight can have several hills, one of them far
from every real basket; training and evaluation that look only near the data never see it. In the mixture
form the weight is the Gaussian density alone, and the posterior of the taste given any basket has exactly
one hill. A worked example with 20 products shows both effects with numbers that can be checked by hand.
\end{abstract}

%======================================================================
\section{Notation}

\begin{itemize}
\item One trip is fixed throughout. $\mathcal J$ is the set of products offered on it. Each product
  $j\in\mathcal J$ belongs to one category $c(j)$, and $\mathcal J_c=\{j: c(j)=c\}$.
\item A \emph{basket} is a non-empty set $S\subseteq\mathcal J$. For a basket,
  $n(S)=|S|$ is its size, $n_c(S)=|S\cap\mathcal J_c|$ the number of its products in category $c$, and
  $m(S)=\#\{c:n_c(S)>0\}$ the number of categories present.
\item $\mathcal S$ is the set of allowed baskets: $1\le n(S)\le 120$ and $n_c(S)\le40$ for every $c$.
  Every sum over baskets below runs over $\mathcal S$.
\item The taste is a vector $z\in\R^r$ ($r=16$ on the real data, $r=1$ in the examples).
  $\N(z;0,I)=(2\pi)^{-r/2}e^{-\lVert z\rVert^2/2}$ is the standard Gaussian density, the ``bell curve''.
\item $\pos{a}=\max(a,0)$ and $\binom{k}{2}=k(k-1)/2$.
\end{itemize}

\paragraph{Parameters.} Each product has a context utility $b_j\in\R$ (price, promotion, season, household
taste, history; a fixed number for this trip) and a taste loading $\phi_j\in\R^r$. Each category has a
taste loading $\psi_c\in\R^r$ for its extra products, a pair penalty $\rho_c$ and count penalties $f_c(k)$
(with $f_c(0)=f_c(1)=0$). $\rho_0(m)$ is the size penalty on the number of categories present.

%======================================================================
\section{The common building block: the basket law given a taste}

\begin{definition}[score and total score]
For a taste $z$ and a basket $S$, the \emph{score} of $S$ is $e^{E_z(S)}$ with
\begin{align}
E_z(S)&=\sum_{j\in S}u_j(z)-\sum_c h_c\big(n_c(S),z\big)-\rho_0\big(m(S)\big),\label{eq:E}\\
u_j(z)&=b_j-\tfrac12\lVert\phi_j\rVert^2+z\cdot\phi_j,\nonumber\\
h_c(k,z)&=\rho_c\tbinom{k}{2}+f_c(k)+\pos{k-1}\big(z\cdot\psi_c-\tfrac12\lVert\psi_c\rVert^2\big).\nonumber
\end{align}
The \emph{total score} of taste $z$ is the sum of the scores of all allowed baskets,
\begin{equation}
Z_z=\sum_{S\in\mathcal S}e^{E_z(S)} .
\label{eq:Z}
\end{equation}
\end{definition}

\begin{definition}[basket law given a taste]
\begin{equation}
p(S\mid z)=\frac{e^{E_z(S)}}{Z_z}.
\label{eq:cond}
\end{equation}
\end{definition}
A basket's probability is its share of the total score; dividing by $Z_z$ makes the probabilities of all
baskets add up to one. Both models below use exactly this law.

\paragraph{The basket field.} Collecting the terms of \eqref{eq:E} that contain $z$ gives
\begin{equation}
E_z(S)=a(S)+z\cdot s(S),\qquad s(S)=\sum_{j\in S}\phi_j-\sum_c\pos{n_c(S)-1}\,\psi_c ,
\label{eq:affine}
\end{equation}
where $a(S)$ collects every term of \eqref{eq:E} without $z$. The vector $s(S)\in\R^r$, the \emph{basket field},
is the direction in taste space that basket $S$ points in.

\subsection{Example 1: three products}

Take $\mathcal J=\{A,B,C\}$, no penalties ($\rho_c=f_c=\rho_0=0$, $\psi_c=0$), and a taste $z$ under which
$u_A(z)=1$, $u_B(z)=0$, $u_C(z)=-1$. Then $e^{E_z(S)}=\exp\sum_{j\in S}u_j(z)$ and
\begin{center}\small
\begin{tabular}{lccccccc|c}\toprule
$S$ & $\{A\}$ & $\{B\}$ & $\{C\}$ & $\{A,B\}$ & $\{A,C\}$ & $\{B,C\}$ & $\{A,B,C\}$ & $Z_z$\\\midrule
$e^{E_z(S)}$ & 2.72 & 1.00 & 0.37 & 2.72 & 1.00 & 0.37 & 1.00 & 9.17\\\bottomrule
\end{tabular}
\end{center}
so $p(\{A\}\mid z)=2.72/9.17=0.30$ and $p(\{A,B,C\}\mid z)=1.00/9.17=0.11$.

Now take a taste $z'$ that raises every utility by 2, so $u_A=3$, $u_B=2$, $u_C=1$:
\begin{center}\small
\begin{tabular}{lccccccc|c}\toprule
$S$ & $\{A\}$ & $\{B\}$ & $\{C\}$ & $\{A,B\}$ & $\{A,C\}$ & $\{B,C\}$ & $\{A,B,C\}$ & $Z_{z'}$\\\midrule
$e^{E_{z'}(S)}$ & 20.1 & 7.4 & 2.7 & 148.4 & 54.6 & 20.1 & 403.4 & 656.7\\\bottomrule
\end{tabular}
\end{center}
Raising each utility by 2 multiplies the score of a basket of $n$ products by $e^{2n}$. The total score grows
about 72-fold, almost all of it from the largest basket, whose probability rises to $403.4/656.7=0.61$.
With thousands of products and baskets of up to 120, a taste that raises many utilities makes $Z_z$
astronomically large, and nearly all of it comes from huge baskets.

\subsection{How $Z_z$ is computed and how it changes with $z$}

$\mathcal S$ is far too large to list, but $Z_z$ is computed exactly. For each category $c$ and each count
$k$, the total score of choosing $k$ of its products is $G_c(k)=e^{-h_c(k,z)}\sum_{A\subseteq\mathcal J_c,|A|=k}\prod_{j\in A}e^{u_j(z)}$,
by a recursion over the products of $c$. The categories are then combined in a table indexed by the
basket size $n$ and the number of categories $m$, and $Z_z=\sum_{n,m}(\text{table entry})\,e^{-\rho_0(m)}$.

For the conditional mean and covariance of the basket field, write
$\E[s\mid z]=\sum_S p(S\mid z)\,s(S)$ and $\Cov[s\mid z]=\sum_S p(S\mid z)\,(s(S)-\E[s\mid z])(s(S)-\E[s\mid z])^\top$.

\begin{lemma}\label{lem:grad}
$\nabla_z\log Z_z=\E[s\mid z]$ and $\nabla_z^2\log Z_z=\Cov[s\mid z]$. In particular $\log Z_z$ is convex in $z$.
\end{lemma}
\begin{proof}
By \eqref{eq:affine}, $\log Z_z=\log\sum_Se^{a(S)+z\cdot s(S)}$. Differentiating once gives
$\sum_S e^{a(S)+z\cdot s(S)}s(S)/Z_z=\E[s\mid z]$; differentiating again gives the covariance, which is positive
semi-definite.
\end{proof}
In words: moving the taste in a direction $v$ raises $\log Z_z$ at the rate $v\cdot\E[s\mid z]$, the expected
basket field along $v$, and that rate itself only increases as $z$ moves further along $v$.

%======================================================================
\section{The old model and its taste weight}

\begin{definition}[old model]
\begin{equation}
p_{\rm old}(S)=\frac{e^{a(S)+\frac12\lVert s(S)\rVert^2}}{Z_{\rm old}},\qquad
Z_{\rm old}=\sum_{S\in\mathcal S}e^{a(S)+\frac12\lVert s(S)\rVert^2}.
\label{eq:old}
\end{equation}
\end{definition}
The term $\tfrac12\lVert s(S)\rVert^2$ is the ``goes-together'' bonus: expanding it, it contains
$\phi_i\cdot\phi_j$ for every pair of products in $S$, so products whose loadings point the same way reward
each other. For a basket of $n$ products with similar loadings it grows like $n^2$, while the penalties in
$a(S)$ grow like $n$.

\begin{lemma}[Gaussian identity]\label{lem:gauss}
For every $s\in\R^r$, $\;e^{\frac12\lVert s\rVert^2}=\int_{\R^r}\N(z;0,I)\,e^{z\cdot s}\,dz$.
\end{lemma}
\begin{proof}
$-\tfrac12\lVert z\rVert^2+z\cdot s=-\tfrac12\lVert z-s\rVert^2+\tfrac12\lVert s\rVert^2$, and $\N(\cdot\,;s,I)$ integrates to one.
\end{proof}

\begin{proposition}[the old model as a weighted mixture]
\begin{equation}
p_{\rm old}(S)=\int_{\R^r}\pi_{\rm old}(z)\,p(S\mid z)\,dz,\qquad
\pi_{\rm old}(z)=\frac{\N(z;0,I)\,Z_z}{Z_{\rm old}} .
\label{eq:oldmix}
\end{equation}
\end{proposition}
\begin{proof}
By Lemma~\ref{lem:gauss} and \eqref{eq:affine},
$e^{a(S)+\frac12\lVert s(S)\rVert^2}=\int\N(z;0,I)\,e^{E_z(S)}\,dz=\int\N(z;0,I)\,Z_z\,p(S\mid z)\,dz$.
Dividing by $Z_{\rm old}$ gives \eqref{eq:oldmix}. Summing over $S$ shows that $\pi_{\rm old}$ integrates to one.
\end{proof}

So the old model also picks a taste and then a basket, but a taste is weighted by the bell curve
\emph{multiplied by its total score} $Z_z$. A taste that makes many high-scoring baskets possible gains weight.

\begin{proposition}[slope and curvature of the old weight]\label{prop:old}
\[
\nabla_z\log\pi_{\rm old}(z)=-z+\E[s\mid z],\qquad
\nabla_z^2\log\pi_{\rm old}(z)=-I+\Cov[s\mid z].
\]
\end{proposition}
\begin{proof}
$\log\pi_{\rm old}(z)=-\tfrac12\lVert z\rVert^2+\log Z_z+\text{const}$; apply Lemma~\ref{lem:grad}.
\end{proof}
Two forces act on the slope: $-z$ pulls back towards zero, and $\E[s\mid z]$ pushes outward, harder the
further out $z$ is (Lemma~\ref{lem:grad}). Where $\Cov[s\mid z]$ exceeds $I$ the curvature is positive, so
$\log\pi_{\rm old}$ is not concave and can have several local maxima, the \emph{hills}. Each hill satisfies
$z=\E[s\mid z]$.

\subsection{Example 2: twenty products and a hidden hill}

Take one taste dimension ($r=1$), 20 products in 20 different categories, all with $b_j=-4$ and $\phi_j=1$,
and no penalties or $\psi$. Then $u_j(z)=-4.5+z$ for every product, $s(S)=n(S)$, and
\[
Z_z=\big(1+e^{-4.5+z}\big)^{20}-1,\qquad
\E[n\mid z]=\frac{20\,\sigma(-4.5+z)\,(1+e^{-4.5+z})^{20}}{Z_z},\qquad\sigma(x)=\frac1{1+e^{-x}},
\]
where the $-1$ removes the empty basket. The slope of Proposition~\ref{prop:old} is $-z+\E[n\mid z]$. Its zeros,
computed numerically:
\begin{center}\small
\begin{tabular}{rrrl}\toprule
$z$ & $\E[n\mid z]$ & $\log\N(z;0,1)Z_z+\tfrac12\log2\pi$ & \\\midrule
1.56 & 1.56 & $-0.63$ & hill 1: near the data\\
2.36 & 2.36 & $-0.68$ & valley\\
20.0 & 20.0 & $+110.0$ & hill 2: all 20 products bought\\\bottomrule
\end{tabular}
\end{center}
Hill 2 is about $e^{110.6}$ times higher than hill 1. Integrating \eqref{eq:oldmix} numerically, the region
$z<5$ holds a fraction $10^{-46}$ of the old model's probability, so under $p_{\rm old}$ essentially every
basket contains all 20 products.

The same follows directly from \eqref{eq:old}. The full basket has $a(S)=20\times(-4.5)=-90$ and bonus
$\tfrac12\cdot20^2=200$, energy $110$; a single product has energy $-4.5+\tfrac12=-4$. The pair bonus wins.

\paragraph{Why fitting and evaluation miss hill 2.} The gradient of the old log-likelihood in any parameter
$\vartheta$ is ``observed minus expected'':
\[
\nabla_\vartheta\log p_{\rm old}(S)=\nabla_\vartheta\big[a(S)+\tfrac12\lVert s(S)\rVert^2\big]-\E_{S'\sim p_{\rm old}}\nabla_\vartheta\big[a(S')+\tfrac12\lVert s(S')\rVert^2\big].
\]
The expectation is over $p_{\rm old}$, and by \eqref{eq:oldmix} over $\pi_{\rm old}$. If real baskets hold one or
two products, chains started from them stay around hill 1, and so does any proposal built from them. The
estimate of the expectation then covers hill 1 only, nothing pushes hill 2 down, and an estimate of
$Z_{\rm old}$ that misses hill 2 makes the reported likelihood too high. On the Dunnhumby refit this hill was
the region of 120-product baskets.

%======================================================================
\section{The mixture model}

\begin{definition}[mixture model]
\begin{equation}
p(S)=\int_{\R^r}\N(z;0,I)\,p(S\mid z)\,dz .
\label{eq:mix}
\end{equation}
\end{definition}
The law given $z$ is the same \eqref{eq:cond}; the taste weight is $\N(z;0,I)$ alone. $Z_z$ appears only as the
divisor inside $p(S\mid z)$, which makes each taste's probabilities add up to one. A taste that allows many
high-scoring baskets gains no weight; it shares its fixed weight among more baskets. $p(S)$ sums to one
because $\sum_Sp(S\mid z)=1$ for every $z$.

\paragraph{Generation.} Draw $z\sim\N(0,I)$, then $S\sim p(\cdot\mid z)$ exactly by the dynamic program of
Section~2.2 (size and categories first, then the count in each category, then the products). No search is
involved, so nothing can be missed.

\paragraph{Example 2 again.} Under \eqref{eq:mix} the taste $z=20$ has density $\propto e^{-200}$. Typical tastes
lie in $[-2,2]$, where each product is bought with probability $\sigma(-4.5+z)\in[0.0015,0.076]$ and
$\E[n\mid z]\in[1.01,1.91]$. The model's mean basket size is $\int\N(z;0,1)\E[n\mid z]\,dz=1.19$, and the
probability of the full 20-product basket is $2.3\times10^{-10}\approx e^{-22.2}$.

\subsection{The posterior of the taste has one hill}

Scoring an observed basket $S$ requires the tastes that explain it. By Bayes' rule their weight is the
posterior
\begin{equation}
p(z\mid S)=\frac{\N(z;0,I)\,p(S\mid z)}{p(S)},\qquad
\log p(z\mid S)=-\tfrac12\lVert z\rVert^2+z\cdot s(S)-\log Z_z+\text{const},
\label{eq:post}
\end{equation}
using \eqref{eq:affine}; the constant does not depend on $z$.

\begin{proposition}[one hill]\label{prop:mix}
\[
\nabla_z\log p(z\mid S)=-z+s(S)-\E[s\mid z],\qquad
\nabla_z^2\log p(z\mid S)=-I-\Cov[s\mid z]\preceq-I .
\]
Hence $\log p(z\mid S)$ is strictly concave, has exactly one maximum $\hat z$, and
$\log p(z\mid S)\le\log p(\hat z\mid S)-\tfrac12\lVert z-\hat z\rVert^2$ for every $z$.
\end{proposition}
\begin{proof}
Differentiate \eqref{eq:post} with Lemma~\ref{lem:grad}. A covariance is positive semi-definite, so the
Hessian is at most $-I$ in every direction. A function whose Hessian is at most $-I$ everywhere has a unique
maximum, and the quadratic bound follows from Taylor's theorem.
\end{proof}

Compare with Proposition~\ref{prop:old}. The term $\E[s\mid z]$ that pushed the old weight outward appears here
with a minus sign: it now works with the bell curve instead of against it. The basket enters only through
the constant vector $s(S)$, which shifts the hill but cannot create a second one. This holds for every basket,
every parameter value and every dimension $r$.

\paragraph{Example 2 again.} For $S=\{\text{one product}\}$ the slope is $-z+1-\E[n\mid z]$. Both $-z$ and
$-\E[n\mid z]$ decrease in $z$, so the slope crosses zero once, at $\hat z=-0.10$.

\subsection{Consequences}

\begin{itemize}
\item \emph{Scoring.} $p(S)$ is estimated by importance sampling around the single hill $\hat z$, found by
  gradient ascent, which always reaches it. There is no second hill to miss.
\item \emph{Training.} The gradient is $\sum_k\bar w_k\big(\nabla E_{z_k}(S)-\E_{S'\sim p(\cdot\mid z_k)}\nabla E_{z_k}(S')\big)$,
  with tastes $z_k$ drawn around $\hat z$ and normalised importance weights $\bar w_k$. The inner expectation is
  exact for each $z_k$ (it is the gradient of $\log Z_{z_k}$), so no region of baskets is left unvisited.
\item \emph{Generation} is exact (above).
\end{itemize}

\begin{center}\small
\begin{tabular}{lll}\toprule
 & old model & mixture model\\\midrule
law given $z$ & $p(S\mid z)$, \eqref{eq:cond} & the same\\
weight of taste $z$ & $\N(z;0,I)\,Z_z/Z_{\rm old}$ & $\N(z;0,I)$\\
slope of the log weight that matters & $-z+\E[s\mid z]$ & $-z+s(S)-\E[s\mid z]$\\
curvature & $-I+\Cov[s\mid z]$ & $-I-\Cov[s\mid z]$\\
number of hills & one or more & exactly one\\
drawing a basket & Markov chain & exact, $z$ then $S\mid z$\\\bottomrule
\end{tabular}
\end{center}

%======================================================================
\section{What changes and what does not}

\begin{itemize}
\item \emph{It is a different model.} The two forms share every parameter but define different laws, so a fit
  of one is not a fit of the other.
\item \emph{Prices and tastes.} In the old model a price change alters $Z_z$ and so shifts which tastes are
  common. In the mixture the taste distribution is fixed by $\N(0,I)$; prices change only what each taste
  buys. This is an assumption, suited to short-run price and promotion questions.
\item \emph{What it does not fix.} The mixture removes hidden hills; it does not by itself make training
  correct. In the full real fit the held-out log-likelihood on all 1{,}024 validation trips fell from
  $-38.10$ at step 240 to $-38.76$ at the end (paired difference $-0.66\pm0.08$), and the drop is real at
  $2{,}048$ importance draws. That is a training problem, to be analysed separately.
\end{itemize}

\end{document}
